\documentclass[11pt]{article}
\usepackage[utf8]{inputenc}
\usepackage[T1]{fontenc}
\usepackage[ngerman]{babel}
\usepackage[margin=2.5cm]{geometry}

\title{Algorithm Engineering\\Woche 2, Aufgabe 4}
\author{Jan Vogt}
\date{}

\begin{document}

\maketitle

\section*{$\pi$ mit \#pragma omp for}

In \texttt{pi\_omp\_for.cpp} ist das Programm \texttt{pi\_numerical\_integration.cpp} aus der
letzten Vorlesung mit \texttt{\#pragma omp for} parallelisiert. Die Schleife wird von OpenMP auf
die Threads aufgeteilt, jeder Thread summiert in eine lokale Variable und am Ende wird diese mit
\texttt{\#pragma omp atomic} zur gemeinsamen Summe addiert:
\begin{verbatim}
#pragma omp parallel num_threads(num_threads)
{
  double sum_local = 0.0;
#pragma omp for schedule(runtime)
  for (int i = 0; i < num_steps; i++) {
    double x = (i + 0.5) * width;
    sum_local = sum_local + (1.0 / (1.0 + x * x));
  }
#pragma omp atomic
  sum += sum_local;
}
\end{verbatim}
Mit \texttt{schedule(runtime)} kann man den Schedule beim Starten über die Umgebungsvariable
festlegen, zum Beispiel \verb|OMP_SCHEDULE=static ./pi_omp_for 8|. Die Threadanzahl wird als
Argument übergeben.

\subsection*{Messung}

Gemessen habe ich auf meinem Laptop (Apple M3) mit 100.000.000 Schritten, ohne Optimierung und
jeweils 5 Läufen. In der Tabelle steht der Median in Sekunden.

\begin{center}
\begin{tabular}{lrrrr}
Schedule & 1 Thread & 2 Threads & 4 Threads & 8 Threads \\
\hline
\texttt{static} & 0.1401 & 0.0709 & 0.0372 & 0.0314 \\
\texttt{dynamic, 1} & 0.5664 & 2.9519 & 2.3110 & 4.1473 \\
\texttt{dynamic, 1000} & 0.1404 & 0.0720 & 0.0371 & 0.0286 \\
\texttt{dynamic, 100000} & 0.1400 & 0.0710 & 0.0365 & 0.0273 \\
\texttt{guided} & 0.1400 & 0.0710 & 0.0368 & 0.0273 \\
\end{tabular}
\end{center}

Mit \texttt{dynamic, 1} ist das Programm mit mehreren Threads viel langsamer als mit einem, weil
sich die Threads für jede einzelne der 100 Millionen Iterationen neue Arbeit holen müssen. Ohne
\texttt{OMP\_SCHEDULE} waren die Zeiten bei mir genauso schlecht wie mit \texttt{dynamic, 1}, man
sollte die Variable bei \texttt{schedule(runtime)} also immer setzen.

Bis 4 Threads sind die anderen Schedules ungefähr gleich schnell, weil jede Iteration gleich viel
Arbeit ist. Mit 8 Threads sind \texttt{dynamic, 100000} und \texttt{guided} aber etwa 13 Prozent
schneller als \texttt{static}. Das ist mir auch im Zusammenhang mit Woche 1 aufgefallen: Dort
mussten bei der statischen Aufteilung vermutlich die Threads auf den Performance-Kernen auf die
langsameren Effizienz-Kerne warten. Mit \texttt{dynamic} können sich die schnellen Kerne einfach
mehr Blöcke holen.
Das passt zu Folie 19, wo steht, dass \texttt{static} bei Effizienz-Kernen nicht unbedingt die
beste Wahl ist.

Außerdem war $\pi$ bei \texttt{static} und gleicher Threadanzahl fast immer exakt gleich, bei
\texttt{dynamic} dagegen in den letzten Stellen bei jedem Lauf etwas anders, weil die Blöcke jedes
Mal anders auf die Threads verteilt werden.

\end{document}
